\section{结论}
\label{sec:conclusion}

本文把 CS149、CS229S 与 CMU 15-849 三门课程的内容统一到「计算、访存、通信」三项分解之下：从 roofline 与 Amdahl 定律出发，经过 SIMD/SIMT 硬件、
GEMM 与 FlashAttention 等 kernel、systolic 加速器、编译器与框架，到大模型的分布式训练、推理服务、模型效率与集群可靠性。
贯穿全文的结论是，几乎每一项系统优化都可以表述为「把一种资源换成另一种」：FlashAttention 用重算换 HBM 流量，ZeRO 用通信换显存，
激活重计算用 FLOPs 换显存，分页 KV 用间接寻址换连续性，推测解码用草稿模型的算力换目标模型的串行步数，MoE 用显存与通信换每 token 的 FLOPs。
判断一次交换是否划算，只需估算交换前后的主导项。

面向面试，本文给出了可当场推导的白板估算、按岗位组织的能力矩阵、五步回答模板与高频错答，并把配套实验整理为可复述的 STAR 故事。
所有配套实验都在 CPU 上完成，它们支撑的是正确性、计数与相对关系的结论，而不是 GPU 上的绝对性能；我们在文中逐处标明了这一区别。
希望这份综述能帮助读者在面试与研究中，从「知道很多名词」走向「能算出哪个名词在当前场景下起作用」。
